\documentclass[11pt]{article}
\usepackage{preamble}
\setlength{\parskip}{4pt}

\begin{document}

\daybanner{AP Statistics \textperiodcentered\ Topic 6.2 (CED 3.3) \textperiodcentered\ B: 2026-11-30 \textperiodcentered\ E: 2027-01-04 \textperiodcentered\ TEACHER KEY}{Narrower Is Not Automatically Better}

\sectionbanner{CED TETHER}

{\small
\begin{itemize}[leftmargin=1.5em,itemsep=2pt,topsep=2pt]
  \item \textbf{Skill 1.D} --- Identify an appropriate inference method for confidence intervals.
  \item \textbf{Skill 4.C} --- Verify that inference procedures apply in a given situation.
  \item \textbf{Skill 3.D} --- Construct a confidence interval, provided conditions for inference are met.
  \item \textbf{EU UNC-4} --- An interval of values should be used to estimate parameters, in order to account for uncertainty.
  \item \textbf{LO UNC-4.A} --- Identify an appropriate confidence interval procedure for a population proportion. [Skill 1.D]
  \item \textbf{EK UNC-4.A.1} --- The appropriate confidence interval procedure for a one-sample proportion for one categorical variable is a one sample z-interval for a proportion.
  \item \textbf{LO UNC-4.B} --- Verify the conditions for calculating confidence intervals for a population proportion. [Skill 4.C]
  \item \textbf{EK UNC-4.B.1} --- In order to make assumptions necessary for inference on population proportions, means, and slopes, we must check for independence in data collection methods and for selection of the appropriate sampling distribution.
  \item \textbf{EK UNC-4.B.2} --- In order to calculate a confidence interval to estimate a population proportion p, we must check for independence and that the sampling distribution is approximately normal: a. To check for independence: (i) Data should be collected using a random sample or a randomized experiment. (ii) When sampling without replacement, check that n $\leq$ 10\%N, where N is the size of the population.
  \item b. To check that the sampling distribution of p-hat is approximately normal (shape): for categorical variables, check that both the number of successes, n$\cdot$p-hat, and the number of failures, n(1 - p-hat) are at least 10.
  \item \textbf{LO UNC-4.C} --- Determine the margin of error for a given sample size and an estimate for the sample size that will result in a given margin of error for a population proportion. [Skill 3.D]
  \item \textbf{EK UNC-4.C.1} --- The standard error of p-hat is SE = sqrt(p-hat(1 - p-hat)/n).
  \item \textbf{EK UNC-4.C.2} --- A margin of error gives how much a value of a sample statistic is likely to vary from the value of the corresponding population parameter.
  \item \textbf{EK UNC-4.C.3} --- For categorical variables, the margin of error is z* $\times$ sqrt(p-hat(1 - p-hat)/n) for a one-sample proportion.
  \item \textbf{EK UNC-4.C.4} --- The formula for margin of error can be rearranged to solve for n, the minimum sample size needed to achieve a given margin of error. Use a guess for p-hat or use p-hat = 0.5 to find an upper bound for the sample size.
  \item \textbf{LO UNC-4.D} --- Calculate an appropriate confidence interval for a population proportion. [Skill 3.D]
  \item \textbf{EK UNC-4.D.1} --- For a one-sample proportion, the interval estimate is p-hat $\pm$ z* $\times$ sqrt(p-hat(1 - p-hat)/n).
  \item \textbf{EK UNC-4.D.2} --- Critical values represent the boundaries encompassing the middle C\% of the standard normal distribution, where C\% is an approximate confidence level for a proportion.
  \item \textbf{LO UNC-4.E} --- Calculate an interval estimate based on a confidence interval for a population proportion. [Skill 3.D]
  \item \textbf{EK UNC-4.E.1} --- Confidence intervals for population proportions can be used to calculate interval estimates with specified units.
\end{itemize}
}
\textbf{Essential question:} Why can a wider interval from a random sample be more trustworthy than a narrower formula result?\par
\textbf{DOK-3 skill:} 4.C \textperiodcentered\ \textbf{FRQ pattern:} {\small\texttt{narrower-\allowbreak{}self-\allowbreak{}selected-\allowbreak{}interval-\allowbreak{}vs-\allowbreak{}wider-\allowbreak{}random-\allowbreak{}interval}}\par
\textbf{DOK rationale:} Part (c) requires judging two competing interval reports by coordinating recruitment, numerical conditions, and interval width, then explaining why the more precise-looking output cannot support the target-population claim.\par

\frameworkphaseheader{Do Now (first take) $\rightarrow$ Explore (video + follow-along) $\rightarrow$ Exit (finish + turn in)}{1 $\rightarrow$ 3}{45}{%
\begin{itemize}[leftmargin=1.5em,itemsep=2pt,topsep=2pt]
  \item Board slide up at the bell. Ask students to choose a plan before calculating.
  \item Begin the 6.1--6.2 video at minute 5; return at minute 33 to separate calculation from generalization.
  \item Collect at the door. Score part (c) only, E/P/I.
\end{itemize}
}{%
\begin{itemize}[leftmargin=1.5em,itemsep=2pt,topsep=2pt]
  \item Choose a plan and cite one collection detail.
  \item Complete the follow-along, finish (a)--(c), revise, and turn in the sheet.
\end{itemize}
}{%
\begin{itemize}[leftmargin=1.5em,itemsep=2pt,topsep=2pt]
  \item Which condition depends on how cardholders entered the sample?
  \item What do the numerical conditions not establish?
  \item Why is Plan B's formula interval narrower?
\end{itemize}
}{Solo teacher. Circulate ~5 pairs. Require all three checks and distinguish a formula output from a defensible interval.}

\begin{callout}{\IconWarn\ WATCH FOR}{calloutred}
\begin{itemize}[leftmargin=1.5em,itemsep=2pt,topsep=2pt]
  \item Treating a larger self-selected group as representative.
  \item Using numerical checks as substitutes for random selection.
  \item Calling Plan B's output a population confidence interval.
\end{itemize}
\end{callout}

\sectionbanner{SCORING PART (c) --- E / P / I}

\textbf{Expected elements}
\begin{itemize}[leftmargin=1.5em,itemsep=2pt,topsep=2pt]
  \item \textbf{selects-plan-a} (required) --- Chooses Plan A for population inference
  \item \textbf{uses-collection-and-numerical-evidence} (required) --- Cites the SRS and one numerical check
  \item \textbf{contrasts-widths-with-validity} (required) --- Notes $0.066<0.137$ but Plan B is nonrandom
  \item \textbf{states-reader-consequence} (required) --- Explains false precision from self-selection
\end{itemize}

{\small\renewcommand{\arraystretch}{1.25}
\noindent\begin{tabular}{|p{0.5in}|p{5.9in}|}
\hline
\textbf{E} & Chooses A, cites both kinds of checks, contrasts widths, and explains self-selection. \\ \hline
\textbf{P} & Chooses A but omits a check, width, or reader consequence. \\ \hline
\textbf{I} & Chooses B for size or width, or treats numerical checks as randomness. \\ \hline
\end{tabular}}

\clearpage
\focusbanner{TODAY'S PROBLEM --- annotated}

A library system has 20,000 adult cardholders and wants to estimate $p$, the proportion who used a digital audiobook last month.\par\smallskip \textbf{Plan A:} A computer takes an SRS of 200 cardholder IDs. All answer; 116 say yes.\par \textbf{Plan B:} A newsletter link lets cardholders choose to respond once. Of 800 respondents, 520 say yes.\par\smallskip \textbf{Eli:} ``Use Plan B. Its larger sample gives the narrower interval.''\par \textbf{Dana:} ``Use Plan A. Entry into the sample matters more than the narrowest output.''\par
\begin{center}
\textbf{Plan counts}\par\smallskip
\begin{tabular}{|l|c|c|c|c|}
\hline
\textbf{Plan} & \textbf{N} & \textbf{n} & \textbf{Yes} & \textbf{No} \\ \hline
\textbf{A} & 20,000 & 200 & 116 & 84 \\ \hline
\textbf{B} & 20,000 & 800 & 520 & 280 \\ \hline
\end{tabular}
\end{center}
\begin{firsttakebox}{FIRST TAKE (5 min)}
Which plan, if either, should estimate all adult cardholders? Cite one collection detail.\par
\answer{Not graded. The larger sample is intentionally attractive.}
\end{firsttakebox}

\textit{Rules callout ``ONE-SAMPLE INTERVAL FOR A PROPORTION (keep on the page)'' is printed on the student sheet.}\par\medskip

\dokbadge{1}~\textbf{(a)} Define the population parameter $p$ in context and name the confidence-interval procedure that would be appropriate when its conditions are met.\par
\answer{Here $p$ is the proportion of all 20,000 adult library cardholders who used a digital audiobook last month. When conditions are met, the procedure is a one-sample $z$-interval for a population proportion.}

\dokbadge{2}~\textbf{(b)} For each plan, compute $\hat p$, check the random, 10 percent, and large-count conditions, and calculate the 95 percent formula interval. Mark any nondefensible population interval.\par
\answer{Plan A has $\hat p=116/200=0.58$. It is an SRS; $200\le0.10(20{,}000)=2{,}000$; and counts are 116 and 84. Thus $SE=0.0348999$, $ME=1.96(SE)=0.0684037$, and the valid interval is $(0.5116,0.6484)$. Plan B has $\hat p=520/800=0.65$; $800\le2{,}000$ and counts 520 and 280 pass, but self-selection fails the random condition. Its mechanical output is $0.65\pm0.0330523$, or $(0.6169,0.6831)$, not a defensible population confidence interval.}

\dokbadge{3}~\textbf{(c)} Decide which plan can support a 95 percent interval for all cardholders. Cite recruitment and a numerical condition, compare interval widths, and explain the reader consequence of reporting Plan B's output.\par
\answer{Use Plan A. Its SRS, $200\le2{,}000$, and counts 116 and 84 support the interval $(0.512,0.648)$ for all cardholders. Its width is about $0.1368$ versus Plan B's smaller formula width, $0.0661$. Plan B's volunteers may differ systematically from nonrespondents; its margin of error covers random variation, not self-selection. Reporting it would falsely suggest that a narrow range represents all cardholders with 95 percent confidence.}

\begin{summaryexitbox}{EXIT}
One thing the video changed about my first take: \hrulefill
\end{summaryexitbox}

\end{document}
